\section{Limitations}
\label{sec:limitations}

What could make this paper wrong? Section~\ref{sec:notclaim} bounds the theorem.
Several of the real risks lie outside it.

\textbf{The headroom band is not yet interpretable.} The oracle-to-baseline
headroom band we have measured so far was produced with a stand-in compiler and a
scripted transport. It is narrow, and it is uninterpretable until the real
pipeline lands end to end. That is the honest statement and we will not sharpen
it. If the band stays narrow once the real compiler and a real model are both in
the loop, no arm between the baseline and the oracle can be distinguished and the
study needs reframing. \pending{headroom band under the integrated pipeline}

\textbf{Two families receive almost no proposals without an alert.} Families 4
and 5 do not touch the demand stream, so the detector fires at its null rate by
construction, every triggered arm degenerates to the baseline controller in those
cells, and pooling them dilutes any effect the alert channel produces elsewhere.
An arrival-side detector on expected against realized arrivals would be legal,
and it is an open scope decision rather than something we have quietly done.

\textbf{Reproducibility, open weights, and the oracle.} The primary hosted
endpoint rejects the decoding seed parameter and a second accepts it without
honoring it, so reproducibility comes from a cache of stored bytes and a reader
who re-derives our numbers from the provider will not get identical text. The
project is cloud-only by a registered, dated operator decision, so the paper
cannot make an open-weight claim. The oracle is an upper bound rather than a
method: it reads the true family and onset, is never tuned on test data, and
refuses the compound family, so compound rows are excluded and we state that
wherever such a row appears. Verification is also late by design, which is why
proposal-to-activation delay is a primary outcome. And the scope is controlled
policy evaluation rather than deployment, since the real trajectories are genuine
sales histories but their disruptions are simulated and true demand stays visible
during a stockout. Redistribution is governed by upstream competition terms, so
an artifact release carries code, generators, prompts, splits, manifests and
hashes and directs users upstream for the data.

\textbf{The registration is not yet frozen, and it gates the confirmatory
claims.} The preregistration file that must carry the endpoints, the stratum
weights, the three contrasts, the Holm family, the budget allocation, the
compiler mapping and the numeric go and kill thresholds does not yet exist. Until
it does the contrast set of Section~\ref{sec:experiments} is provisional, no
analysis can legitimately be labeled confirmatory, and the per-method freeze
manifest is likewise unimplemented. The verifier and the schema parser are
implemented on feature branches and are not yet integrated on the main line, and
no pilot has been run, so no go or kill verdict exists. We say this plainly
because the alternative is a paper whose method section describes a system its
own evidence base does not yet contain. \pending{integration status, pilot
verdict, frozen registration hash}
